\chapter{A17 智慧= 16 x 理解} % Introduction chapter suppressed from the table of contents

\hypertarget{ux5bf9ux7ba1ux7406ux8005ux6765ux8bf4ux4e00ux76ceux53f8ux7684ux667aux6167ux7b49ux4e8eux4e00ux78c5ux7684ux7406ux89e3}{%
\section{对管理者来说，一盎司的智慧等于一磅的理解}\label{ux5bf9ux7ba1ux7406ux8005ux6765ux8bf4ux4e00ux76ceux53f8ux7684ux667aux6167ux7b49ux4e8eux4e00ux78c5ux7684ux7406ux89e3}}

Ackoff教授（第73条）：To managers, an ounce of wisdom is worth a pound
of understanding.

\framebox{%
\begin{minipage}[t]{0.97\columnwidth}\raggedright
Ackoff 教授经验之谈：\\
使用以下公式:

\begin{itemize}
\tightlist
\item
  1盎司的智慧 = 1磅的理解
\item
  1盎司的理解 = 1磅的知识
\item
  1盎司知识 = 1磅信息
\item
  1盎司信息 = 1磅数据
\end{itemize}

\begin{description}
\tightlist
\item[]
能算出一盎司的智慧 = 65,536盎司的数据。
\end{description}

智慧包含在价值陈述中，例如格言和谚语：

\begin{itemize}
\tightlist
\item
  它使我们能够感知和评估所做事情的长期和短期后果
\item
  它使我们追求具有持久价值的东西
\item
  让我们为长期收益作出短期牺牲，避免因贪图眼前利益而牺牲未来
\end{itemize}

知识帮助我们能高效做事。\\
理解让我们可以按自己的想法与愿望把事情做好;\\
智慧则指引我们追求``正确''的事情，从而提升能获得我们所需的能力。\\
信息、知识和理解让我们能高效做事，而智慧则让我们做正确的事情，并有效。\\
科学追求数据、信息、知识和理解，探求真相，人文学则追求智慧,探索真理。
\strut
\end{minipage}}

\newpage

每当听到家长说：``现代科技日新月异，小孩应专注于学好数理化生，毕业后才能在社会中立足''。

我就会立马想到史学家钱穆先生对中西方教育的分析：西方教育如同大型超市，提供各样学科，满足学生需求。而中国教育几千年来更强调的是如何做人，老师不仅传授知识与技术，还要以身作则，做好榜样，期望学生日后能回馈家庭与社会。

在现代西方教育体系中，若缺乏家庭教育，学生就可能善恶不分，即使掌握了尖端科技与知识，也可能危害社会。

``学了教科书里的知识并能在考试中写出正确答案，就算理解了。''
你赞同吗？\\
二十多年前，我在2018年靠开始讲ACP敏捷课程，之后一直做培训与辅导，帮开发团队从传统瀑布式开发转变为敏捷开发，因为要面对并解答各类疑问，我不断自学其他如精益、XY理论、必须改善整个系统等管理思想，加深了我对敏捷项目管理的理解，例如敏捷宣言第二条：``工作的软件高于详尽的文档''\\
最初，我浅显地认为代码比文档更重要。随着积累的培训越多，我就越能体会到这条原则体现出的精益思想。80年代，很多软件开发项目花了大量精力做设计，画出很多``波波''设计图，但没有写代码，导致最终项目延期或失败。按精益思想，应通过小步验证可运行得代码，不能单靠需求或设计评审，才能避免开发完成后产生大量缺陷和本可避免的返工。\\
第三条``客户合作高于合同谈判''\\
初学敏捷时，我以为因需求对开发至关重要，所以不能单靠需求文档，还需与客户沟通以防误解。后来，我认识到系统整体最优的理念：即使各部分单独达标或做到最优，但若缺乏协调，整体未必最优。\\
客户、产品经理、设计、编码等团队间必须通力合作，才能实现整体最优，而非仅靠个体达标。
这与XY理论强调的团队协作一致，也呼应了第一条``个体和互动高于流程和工具''。敏捷大师们从以往开发项目失败的经验，知道因技术与需求有很多变数，不可能制定3-6个月详细计划，应按精益思想，先针对本次两周冲刺制定详细计划，并让团队尽力在每次冲刺开发出有价值的产品。但有些人却误以为敏捷开发项目都不需要制定计划。\\
所以要了解敏捷宣言背后的思想必须先了解精益、XY和系统总体最优等。学习越深入，越发现这些概念是息息相关的。理解了背后原因与核心思想，才是从理解迈向智慧的关键。从理解中提炼的智慧，真正有用的不足十分之一。


